\documentclass[UTF8,zihao=-4,openany]{ctexbook}

\usepackage[a4paper,margin=2.45cm,headheight=15pt]{geometry}
\usepackage{amsmath,amssymb,bm,mathtools}
\usepackage{graphicx}
\usepackage{booktabs,tabularx,array}
\graphicspath{{figures/}}
\newcommand{\bookfig}[4]{%
  \begin{samepage}
  \begin{center}
    \includegraphics[width=0.92\linewidth,height=0.38\textheight,keepaspectratio]{#1}%
  \end{center}
  \noindent{\heiti 图~#2}\quad #3\par
  \vspace{0.3em}
  \noindent 符号：#4\par
  \vspace{0.5em}%
  \end{samepage}%
}
\usepackage[dvipsnames]{xcolor}
\usepackage{fancyhdr}
\usepackage{enumitem}
\usepackage{listings}
\usepackage[
  unicode,
  colorlinks=true,
  linkcolor=Primary,
  urlcolor=Accent,
  bookmarksopen=true,
  bookmarksnumbered=true,
  bookmarksdepth=1
]{hyperref}

\definecolor{Primary}{HTML}{355C7D}
\definecolor{Accent}{HTML}{6C8EAD}
\definecolor{CodeBg}{HTML}{F6F7F8}

\setlength{\parindent}{2em}
\setlength{\parskip}{0.35em}
\setlist[itemize]{leftmargin=2em,itemsep=0.2em,topsep=0.3em}
\setlist[enumerate]{leftmargin=2.2em,itemsep=0.2em,topsep=0.3em}
\renewcommand{\arraystretch}{1.25}
\setcounter{secnumdepth}{2}
\setcounter{tocdepth}{1}
\allowdisplaybreaks[4]
\emergencystretch=3em
\hfuzz=5pt
\vfuzz=5pt

\pagestyle{fancy}
\fancyhf{}
\fancyhead[L]{数值线性代数}
\fancyhead[R]{\hyperlink{contents}{返回目录}}
\fancyfoot[C]{\thepage}

\hypersetup{
  pdftitle={数值线性代数学习笔记},
  pdfauthor={learning-notes}
}

\lstset{
  basicstyle=\ttfamily\small,
  backgroundcolor=\color{CodeBg},
  frame=single,
  rulecolor=\color{Accent},
  breaklines=true,
  columns=fullflexible,
  keepspaces=true,
  showspaces=false,
  showstringspaces=false
}

\newcommand{\code}[1]{\texttt{#1}}
\newcommand{\pending}{%
本节标题依徐树方、高立、张平文\emph{数值线性代数}第二版目录保留。
正文待后续补充，此处不编造结论。}

\title{\heiti 数值线性代数\\[0.4em]\Large 学习笔记}
\author{内容提炼}
\date{\today}

\begin{document}

\maketitle
\frontmatter

\chapter*{说明}
\addcontentsline{toc}{chapter}{说明}
本笔记对应徐树方、高立、张平文编著
\emph{数值线性代数}（第二版，北京大学出版社，2011~年修订；
北京大学数学教学系列丛书）。
目录层级与顺序严格按原书印刷目录。
试作版充实绪论与第~1.1~节（三角形方程组与三角分解/LU）；
第~1.2~节起及第~2--7~章仅保留原书标题。
正文以公式陈述定义、成立条件、关键推导与结论；
每一公式后写明输入如何算出输出，以及该数在问题中的含义。

\hypertarget{contents}{}
\tableofcontents

\mainmatter
\chapter*{绪论}
\addcontentsline{toc}{chapter}{绪论}
\markboth{绪论}{绪论}
\setcounter{section}{0}
\renewcommand{\thesection}{\arabic{section}}

数值线性代数又称矩阵计算（matrix computation）。
中心任务不是再证一遍存在唯一性，而是：针对矩阵问题的结构，设计\textbf{快速且可靠}的算法。

\section{数值线性代数的基本问题}

三类核心问题：
\begin{align}
Ax&=b,\\
\min_y\,\lVert Ay-b\rVert_2,\\
Ax&=\lambda x\quad(x\neq 0).
\end{align}
第一式输入非奇异 $A\in\mathbf{R}^{n\times n}$ 与 $b\in\mathbf{R}^{n}$，输出未知量 $x$。
第二式输入 $A\in\mathbf{R}^{m\times n}$、$b\in\mathbf{R}^{m}$，输出使残差欧氏长度最小的 $x$；该长度非负，一般不能把残差降到零。
第三式输出特征值 $\lambda$ 与特征向量 $x$（可差非零倍数）。
奇异值分解因应用广泛，现常列为第四类问题。

\section{研究数值方法的必要性}

理论上的闭式解往往不能当算法用。
Cramer 法则把 $x_i$ 写成两个行列式之比，再配 Laplace 展开，解 $n$ 阶组的乘法次数大于 $(n+1)!$。
在每秒 $10^{10}$ 次运算的机器上取 $n=25$，约需 $13$ 亿年；同一问题用消元法不到一秒。
$n!$ 是运算次数，不是解的某个分量——它说明「公式漂亮」与「可计算」不是一回事。
特征值的 Jordan 分解能读出重数，但对扰动不稳定；实际计算改用 Schur 分解。

\section{矩阵分解是设计算法的主要技巧}

全书的设计思想可以压成一句：
把一般矩阵问题化成一个或几个\textbf{容易求解的特殊问题}，而完成这一转化的主要技巧是矩阵分解。

\bookfig{fig-intro-triangular.png}{绪论第4页}{下三角组与上三角组可以前代、回代直接解出。}{$L$：下三角；$U$：上三角；$x,y$：未知量；$b,c$：右端}

一般的 $Ax=b$ 先分解
\[
PA=LU,
\]
再解 $Ly=Pb$ 与 $Ux=y$。
$P$ 是排列方阵，$L$ 下三角，$U$ 上三角；$y$ 只是中间量，$x$ 才是原未知量。
于是「解线性组」变成「如何分解」加上两次 $O(n^2)$ 的三角形求解。这是第一章的主线。

\section{敏度分析与误差分析}

计算结果很少等于真值。误差来自数据本身，也来自计算过程。
要问「差多少」，必须把\textbf{问题}和\textbf{算法}分开看。

敏度分析：数据微扰后，真解变多少。在相对扰动很小的前提下，
\begin{equation}
\frac{\lvert f(x+\delta x)-f(x)\rvert}{\lvert f(x)\rvert}
\le c(x)\frac{\lvert\delta x\rvert}{\lvert x\rvert}.
\end{equation}
$c(x)$ 是 $f$ 在 $x$ 处的条件数（condition number），无量纲。
$c(x)$ 大称病态，小称良态——这是问题的固有属性，与算法无关。
可微时 $c(x)\approx\lvert f'(x)\rvert\,\lvert x\rvert/\lvert f(x)\rvert$。

误差分析（后向）：把算出的 $\hat y$ 看成精确作用在受扰数据上的结果，
\[
\hat y=f(x+\delta x),\qquad \lvert\delta x\rvert\le\lvert x\rvert\,\varepsilon.
\]
$\varepsilon$ 小称算法数值稳定——这是算法的固有属性，与问题是否病态无关。

合在一起：
\begin{equation}
\frac{\lvert\hat y-f(x)\rvert}{\lvert f(x)\rvert}\le c(x)\,\varepsilon.
\end{equation}
相对误差上界是条件数与后向扰动水平之积。
只有对良态问题使用稳定算法，才可期望可靠结果。

\section{算法复杂性与收敛速度}

直接法：无舍入时有限步得到精确解，快慢主要看运算量（现计全部四则，并略去低阶，如 $3n^3+20n^2$ 称作 $3n^3$）。
运算量不能单独决定快慢：数据传输往往比算术更贵。

迭代法：有限步得不到精确解，还要看
\[
\lVert x_k-x\rVert\le c\lVert x_{k-1}-x\rVert
\quad(0<c<1)
\]
（线性收敛）还是
\[
\lVert x_k-x\rVert\le c\lVert x_{k-1}-x\rVert^2
\]
（平方/二次收敛）。$\lVert x_k-x\rVert$ 是第 $k$ 步误差的长度；$c$ 越小越好。

\section{算法的软件实现与现行数值线性代数软件包}

可靠实现需要算法理论、机器结构与计算经验，不是把公式翻译成循环。
基本算法主要收在 LAPACK 与 MATLAB 中；后者便于试用，一般慢于直接调 LAPACK。

\section{符号说明}

矩阵用大写 $A,B,\Lambda$，元素 $a_{ij}$，向量 $x,y$，单位阵 $I$，其第 $i$ 列为 $e_i$。
排列方阵 $P=[e_{i_1},\ldots,e_{i_n}]$。
$\operatorname{tr}A$、$\operatorname{rank}A$、$\det A$、$A^{-1}$、$A^{\mathsf T}$、$A^{-\mathsf T}=(A^{-1})^{\mathsf T}$。
算法下标沿用 $i:j$、$A(i,:)$、$A(k:l,p:q)$，分别表示整数段、一行、一个子块。

\renewcommand{\thesection}{\thechapter.\arabic{section}}

\chapter{线性方程组的直接解法}

直接法指无舍入时有限步得到 $Ax=b$ 的精确解。
本章的对象是中小规模、系数无特殊结构的方程组，方法是 Gauss 消元。

\section{三角形方程组和三角分解}

中心思想：三角形组 $O(n^2)$ 可解，所以先把 $A$ 变成 $LU$；
工具是 Gauss 变换——用一个乘数向量把一列下方打成零，而其逆几乎不花代价。
$L_k$ 与 $A^{(k)}$ 都不需要按满阵去乘，由同一条秩~$1$ 公式一次写出。

\subsection{三角形方程组的解法}

方阵非奇异（nonsingular）指可逆：存在 $B$ 使 $AB=BA=I$，等价于 $\det A\neq 0$。
下三角阵 $L$ 非奇异当且仅当对角元 $l_{ii}\neq 0$（$i=1,\ldots,n$）。

非奇异下三角组
\begin{equation}
Ly=b
\end{equation}
中 $L=[l_{ij}]\in\mathbf{R}^{n\times n}$ 与 $b\in\mathbf{R}^{n}$ 已知，$y\in\mathbf{R}^{n}$ 未知。
由第 $i$ 个方程解出
\begin{equation}
y_i=\Bigl(b_i-\sum_{j=1}^{i-1}l_{ij}y_j\Bigr)\big/l_{ii},\qquad i=1,\ldots,n.
\end{equation}
输入已求得的 $y_1,\ldots,y_{i-1}$ 与第 $i$ 行，输出标量 $y_i$。
这就是前代法（forward substitution）。

\bookfig{fig-11-lower-L.png}{第11页}{下三角组 $Ly=b$：$L$ 已知且对角非零，$y$ 未知，$b$ 已知。}{$L$：非奇异下三角；$y$：未知量；$b$：右端}

将 $y_i$ 就地写回 $b_i$，即原书算法~1.1.1。运算量
\[
\sum_{i=1}^{n}(2i-1)=n^{2}.
\]
$n^2$ 是四则次数，不是解的某个分量。

\bookfig{fig-11-alg-111.png}{第12页}{前代：用第 $j$ 个已求出的分量去改它下面尚未求解的分量。}{$b$：右端（就地覆盖为 $y$）；$L(j,j)$：对角主元}

非奇异上三角组 $Ux=y$（$u_{ii}\neq 0$）由最后一个方程倒着解，称为回代法（back substitution）：
\begin{equation}
x_i=\Bigl(y_i-\sum_{j=i+1}^{n}u_{ij}x_j\Bigr)\big/u_{ii},\qquad i=n,n-1,\ldots,1.
\end{equation}
输入已求得的 $x_{i+1},\ldots,x_n$ 与第 $i$ 行，输出标量 $x_i$。运算量同样是 $n^2$。

\bookfig{fig-11-alg-112.png}{第13页}{回代：从最后一个未知量倒着做，运算量同样是 $n^2$。}{$y$：右端（就地覆盖为 $x$）；$U(j,j)$：对角主元}

于是，对一般的
\begin{equation}
Ax=b,
\end{equation}
若能写成 $A=LU$（$L$ 下三角、$U$ 上三角），则先解 $Ly=b$ 得中间量 $y$，再解 $Ux=y$ 得未知量 $x$。
两次三角形求解共 $2n^2$ 次运算。本节剩下的事只有一件：如何构造 $L$ 与 $U$。

\subsection{Gauss 变换}

要把 $A$ 逐步变成上三角，同时让变换本身保持下三角。
对给定 $x=(x_1,\ldots,x_n)^{\mathsf T}$，用最简单的单位下三角把第 $k+1$ 至 $n$ 个分量变成 $0$。
这样的矩阵是
\begin{equation}
L_k=I-l_k e_k^{\mathsf T},
\end{equation}
其中 Gauss 向量
\[
l_k=(0,\ldots,0,l_{k+1,k},\ldots,l_{nk})^{\mathsf T}.
\]

\bookfig{fig-11-gauss-Lk.png}{第14页}{Gauss 变换：第 $k$ 列对角下方是乘数的相反数，其余同单位阵。}{$L_k$：Gauss 变换；$l_{ik}$：乘数}

成立条件 $x_k\neq 0$。取无量纲比值
\begin{equation}
l_{ik}=\frac{x_i}{x_k},\qquad i=k+1,\ldots,n,
\end{equation}
则
\begin{equation}
L_k x=(x_1,\ldots,x_k,0,\ldots,0)^{\mathsf T}.
\end{equation}
输入 $x$ 与主元 $x_k$，输出前 $k$ 个分量不变、其后全为 $0$ 的向量。

因为 $e_k^{\mathsf T}l_k=0$，
\[
(I-l_k e_k^{\mathsf T})(I+l_k e_k^{\mathsf T})=I,
\]
故
\begin{equation}
L_k^{-1}=I+l_k e_k^{\mathsf T}.
\end{equation}
逆就是把乘数改号，不必另做求逆。

作用在矩阵上是秩~$1$ 修正：
\begin{equation}
L_k A=A-l_k(e_k^{\mathsf T}A).
\end{equation}
$e_k^{\mathsf T}A$ 是 $A$ 的第 $k$ 行（$1\times n$）；$l_k$ 与之相乘得到 $n\times n$ 的秩~$1$ 矩阵。
这就是用第 $k$ 行去改它下面的行，不必把 $L_k$ 当成稠密阵去乘。

\subsection{三角分解的计算}

$A=LU$ 称为三角分解，亦称 LU 分解，其中 $L$ 为下三角、$U$ 为上三角。
记 $A^{(0)}=A$，目标是求 $n-1$ 个 Gauss 变换，使
\begin{equation}
A^{(k)}=L_k A^{(k-1)}=L_k\cdots L_1 A
\end{equation}
逐步变成上三角。

\bookfig{fig-11-block-Ak.png}{第16页}{前 $k-1$ 步之后，左上已是上三角，右下 $A_{22}^{(k-1)}$ 是待消的块。}{$A^{(k-1)}$：当前矩阵；$A_{11}^{(k-1)}$：已完成的上三角；$A_{22}^{(k-1)}$：尚未消去的主子块}

\textbf{快速求 $L_k$。}
成立条件：主元 $a_{kk}^{(k-1)}\neq 0$。输入 $A^{(k-1)}$ 的第 $k$ 列，输出 $n-k$ 个乘数
\begin{equation}
l_{ik}=\frac{a_{ik}^{(k-1)}}{a_{kk}^{(k-1)}},\qquad i=k+1,\ldots,n,
\end{equation}
从而 $L_k=I-l_k e_k^{\mathsf T}$ 唯一确定。
独立数据只有这 $n-k$ 个数，不必生成满的 $n\times n$ 阵。

\textbf{快速求 $A^{(k)}$。}
由秩~$1$ 公式，前 $k$ 行不变；第 $k$ 列下方变成 $0$；尾块按元素更新为
\begin{equation}
a_{ij}^{(k)}=a_{ij}^{(k-1)}-l_{ik}\,a_{kj}^{(k-1)},
\qquad i,j=k+1,\ldots,n.
\end{equation}
输入 $A^{(k-1)}$ 的第 $k$ 行与乘数 $l_k$，输出阶为 $n-k$ 的新尾块 $A_{22}^{(k)}$。
第 $k$ 步运算量是 $(n-k)+2(n-k)^2$，不是一次 $O(n^3)$ 的稠密乘法。

\bookfig{fig-11-block-Ak-step.png}{第16页}{再乘 $L_k$ 后，上三角块扩大一阶。}{$A^{(k)}$：第 $k$ 步结果；$k$ 与 $n-k$：块的阶}

若前 $n-1$ 个主元全非零，则 $U=A^{(n-1)}$ 为上三角，且
\begin{equation}
A=L_1^{-1}\cdots L_{n-1}^{-1}U=LU,
\end{equation}
其中
\[
L=L_1^{-1}\cdots L_{n-1}^{-1}.
\]
各 $L_k^{-1}$ 只是把乘数写回对角下方；因 $j<i$ 时 $e_j^{\mathsf T}l_i=0$，乘积仍是单位下三角，乘数按列填进 $L$：
\begin{equation}
L=I+[l_1,\ldots,l_{n-1},0].
\end{equation}

\bookfig{fig-11-unit-L.png}{第17页}{$L$ 的亚对角就是各步 Gauss 向量，对角全为 $1$。}{$l_i$：第 $i$ 步 Gauss 向量；$L$：单位下三角}

$A^{(k)}$ 第 $k$ 列下方已是 $0$，无需再存；乘数就覆盖在这些位置。一份数组里同时装着 $L$ 与 $U$。

\bookfig{fig-11-inplace-lu.png}{第18页}{就地存储：严格下三角位置是 $L$ 的乘数，对角及上方是 $U$；算法~1.1.3 按列写入。}{$l_{ij}$：乘数；$a_{ij}^{(k)}$：第 $k$ 步后仍留在 $U$ 位置上的元}

算法~1.1.3 的运算量
\begin{equation}
\sum_{k=1}^{n-1}\bigl((n-k)+2(n-k)^{2}\bigr)
=\frac{n(n-1)}{2}+\frac{n(n-1)(2n-1)}{3}
=\frac{2}{3}n^{3}+O(n^{2}).
\end{equation}
这与 Cramer 法则的 $(n+1)!$ 不在同一量级。
分解一次后，每个新右端只需再花 $2n^{2}$ 做前代与回代。

主元全非零，算法才能进行到底。记 $A_k$ 为 $A$ 的 $k$ 阶顺序主子阵（leading principal submatrix），即左上 $k\times k$ 子阵。

\noindent{\heiti 定理 1.1.1}\quad
主元 $a_{ii}^{(i-1)}$（$i=1,\ldots,k$）全非零，当且仅当顺序主子阵 $A_1,\ldots,A_k$ 都非奇异。

\bookfig{fig-11-leading-minor.png}{第19页}{做到第 $k-1$ 步时，$k$ 阶顺序主子阵是「已完成的上三角加上当前主元」；左乘各 $L_j$ 的主子阵不改变这一形状。}{$A_{11}^{(k-1)}$：前 $k-1$ 个主元所在的上三角；$a_{kk}^{(k-1)}$：当前主元；$A_k$：$k$ 阶顺序主子阵}

因为单位下三角的行列式为 $1$，
\begin{equation}
\det A_k=a_{kk}^{(k-1)}\det A_{11}^{(k-1)}.
\end{equation}
$\det A_k$ 是 $k$ 阶子阵是否可逆的判定量；当前主元是否为零，就是这块是否非奇异。

\noindent{\heiti 定理 1.1.2}\quad
若 $A_k$（$k=1,\ldots,n-1$）均非奇异，则存在唯一的单位下三角 $L$ 与上三角 $U$，使得 $A=LU$。

存在性就是上面的消元。唯一性：若 $A=LU=\tilde L\tilde U$，则
\[
\tilde L^{-1}L=\tilde U U^{-1}.
\]
左边是单位下三角，右边是上三角，故两边只能是 $I$，从而 $L=\tilde L$、$U=\tilde U$。

$A$ 本身非奇异并不蕴含顺序主子阵都非奇异，因此可逆也不能保证无选主元的消元能做完。
下一节处理小主元。

\section{选主元三角分解}

\pending

\section{平方根法}

\pending

\section{分块三角分解}

\pending

\section{习题}

\pending

\section{上机习题}

\pending

\chapter{线性方程组的敏度分析与消去法的舍入误差分析}

\section{向量范数和矩阵范数}

\subsection{向量范数}

\pending

\subsection{矩阵范数}

\pending

\section{线性方程组的敏度分析}

\pending

\section{基本运算的舍入误差分析}

\pending

\section{列主元 Gauss 消去法的舍入误差分析}

\pending

\section{计算解的精度估计和迭代改进}

\subsection{精度估计}

\pending

\subsection{迭代改进}

\pending

\section{习题}

\pending

\section{上机习题}

\pending

\chapter{最小二乘问题的解法}

\section{最小二乘问题}

\pending

\section{初等正交变换}

\subsection{Householder 变换}

\pending

\subsection{Givens 变换}

\pending

\section{正交变换法}

\pending

\section{习题}

\pending

\section{上机习题}

\pending

\chapter{线性方程组的古典迭代解法}

直接法有限步得到精确解，运算量由 \(n^3\) 主导。
大规模稀疏组若仍做 \(LU\)，填入会毁掉稀疏性。
古典迭代把 \(Ax=b\) 改成
\[
x^{(k+1)}=Gx^{(k)}+d,
\]
每步只做矩阵–向量或三角回代，有限步一般得不到真解 \(x^{*}\)，停机看残差或相邻两步差。

\section{单步线性定常迭代法}

把 \(A\) 写成
\begin{equation}
A=D-L-U.
\end{equation}
\(D=\operatorname{diag}(a_{11},\ldots,a_{nn})\) 为对角；
\(L\) 取 \(A\) 严格下三角部分的相反数，\(U\) 取严格上三角部分的相反数。
要求 \(a_{ii}\neq 0\)。\(L,U\) 的元是非负的当且仅当 \(A\) 的非对角元非正。

\subsection{Jacobi 迭代法}

取 \(M=D\)，则
\begin{equation}
x^{(k+1)}=D^{-1}(L+U)x^{(k)}+D^{-1}b.
\end{equation}
迭代矩阵 \(G_J=D^{-1}(L+U)=I-D^{-1}A\)。
第 \(i\) 个分量只用上一步的全部未知量：
\begin{equation}
x_i^{(k+1)}=\frac{1}{a_{ii}}\Bigl(b_i-\sum_{j\neq i}a_{ij}x_j^{(k)}\Bigr).
\end{equation}
输入 \(A\) 的第 \(i\) 行、\(b_i\) 与 \(x^{(k)}\)，输出标量 \(x_i^{(k+1)}\)。各分量彼此独立，可并行。

\subsection{Gauss-Seidel 迭代法}

取 \(M=D-L\)，则
\begin{equation}
x^{(k+1)}=(D-L)^{-1}Ux^{(k)}+(D-L)^{-1}b.
\end{equation}
迭代矩阵 \(G_{\mathrm{GS}}=(D-L)^{-1}U\)。
算第 \(i\) 个分量时，下标小于 \(i\) 的已换成本步新值：
\begin{equation}
x_i^{(k+1)}=\frac{1}{a_{ii}}\Bigl(b_i-\sum_{j<i}a_{ij}x_j^{(k+1)}-\sum_{j>i}a_{ij}x_j^{(k)}\Bigr).
\end{equation}
输入本步已更新的 \(x_1^{(k+1)},\ldots,x_{i-1}^{(k+1)}\) 与尚未更新的 \(x_{i+1}^{(k)},\ldots,x_n^{(k)}\)，输出 \(x_i^{(k+1)}\)。
\(D-L\) 是下三角，每步一次前代，不能像 Jacobi 那样整向量同时更新。

\subsection{单步线性定常迭代法}

一般格式是
\begin{equation}
x^{(k+1)}=Gx^{(k)}+d,
\end{equation}
其中 \(G\in\mathbf{R}^{n\times n}\) 不随 \(k\) 改变（定常），\(d\) 为常向量。
由分裂 \(A=M-N\) 得到 \(G=M^{-1}N\)、\(d=M^{-1}b\)。
若极限存在，记 \(x^{*}=\lim x^{(k)}\)，则 \(x^{*}=Gx^{*}+d\)。
当 \(I-G\) 可逆时，该极限就是 \(Ax=b\) 的解。
Jacobi、G--S 都是这一格式的特例。

\section{收敛性理论}

\subsection{收敛的充分必要条件}

误差 \(e^{(k)}=x^{(k)}-x^{*}\) 满足 \(e^{(k)}=G^k e^{(0)}\)。
对任意初值 \(x^{(0)}\) 都有 \(x^{(k)}\to x^{*}\)，当且仅当谱半径
\begin{equation}
\rho(G)=\max_i\lvert\lambda_i(G)\rvert<1.
\end{equation}
\(\rho(G)\) 是 \(G\) 特征值模的最大值，无量纲。\(\rho(G)\ge 1\) 时至少存在一个初值使迭代不收敛。

\subsection{收敛的充分条件及误差估计}

任一矩阵范数下 \(\rho(G)\le\lVert G\rVert\)。因此
\begin{equation}
\lVert G\rVert<1
\end{equation}
是收敛的充分条件（不是必要）。此时
\begin{equation}
\lVert e^{(k)}\rVert\le\lVert G\rVert^k\lVert e^{(0)}\rVert,
\qquad
\lVert e^{(k)}\rVert\le\frac{\lVert G\rVert}{1-\lVert G\rVert}\lVert x^{(k)}-x^{(k-1)}\rVert.
\end{equation}
前式用 \(k\) 步收缩比估计与真解的距离；后式用相邻两步差估计，便于停机。

\subsection{Jacobi 迭代与 G-S 迭代法的收敛性}

若 \(A\) 严格对角占优，即
\[
\lvert a_{ii}\rvert>\sum_{j\neq i}\lvert a_{ij}\rvert,\qquad i=1,\ldots,n,
\]
则 \(\lVert G_J\rVert_{\infty}<1\)、\(\lVert G_{\mathrm{GS}}\rVert_{\infty}<1\)，两种迭代都收敛。

若 \(A\) 对称正定，则 G--S 收敛（Jacobi 还要额外条件，例如 \(2D-A\) 正定）。

若 \(a_{ii}>0\) 且 \(a_{ij}\le 0\)（\(i\neq j\)），则 Jacobi 与 G--S 同收敛或同发散；收敛时 \(\rho(G_{\mathrm{GS}})\le\rho(G_J)<1\)，G--S 不慢于 Jacobi。

\section{收敛速度}

\subsection{平均收敛速度和渐近收敛速度}

\(k\) 步平均收敛速度
\begin{equation}
R(G^k)=-\ln\bigl(\lVert G^k\rVert^{1/k}\bigr).
\end{equation}
渐近收敛速度
\begin{equation}
R_{\infty}(G)=-\ln\rho(G).
\end{equation}
两者都是正数（当 \(\rho(G)<1\)）；越大越快。主导的是 \(R_{\infty}\)：\(\rho(G)\) 越接近 \(1\)，要达到给定精度所需步数越多。

\subsection{模型问题}

单位正方形上 Poisson 方程的五点差分：步长 \(h=1/(m+1)\)，未知量个数 \(n=m^{2}\)。
系数阵对称正定、稀疏、严格对角占优。这一模型用来比较各迭代法的 \(\rho(G)\) 随 \(h\) 怎样变坏，不是另给一套算法。

\subsection{Jacobi 迭代法和 G-S 迭代法的渐近收敛速度}

对该模型，
\begin{equation}
\rho(G_J)=\cos(\pi h),\qquad
\rho(G_{\mathrm{GS}})=\cos^{2}(\pi h)=\rho(G_J)^{2}.
\end{equation}
\(h\to 0\) 时两者的谱半径都 \(\to 1\)，速度都变慢；同一网格上 G--S 的渐近速度大约是 Jacobi 的两倍。

\section{超松弛迭代法}

\subsection{迭代格式}

在 G--S 新值与旧值之间作加权，松弛因子为 \(\omega\)：
\begin{equation}
(D-\omega L)x^{(k+1)}=\bigl[(1-\omega)D+\omega U\bigr]x^{(k)}+\omega b.
\end{equation}
分量形式为
\begin{equation}
x_i^{(k+1)}=x_i^{(k)}+\frac{\omega}{a_{ii}}\Bigl(b_i-\sum_{j<i}a_{ij}x_j^{(k+1)}-\sum_{j\ge i}a_{ij}x_j^{(k)}\Bigr).
\end{equation}
\(\omega=1\) 就是 G--S；\(0<\omega<1\) 称低松弛，\(1<\omega<2\) 称超松弛（SOR）。
迭代矩阵记为 \(\mathcal L_{\omega}\)。

\subsection{收敛性分析}

若 \(A\) 对称正定，则 SOR 对任意初值收敛的充要条件是
\begin{equation}
0<\omega<2.
\end{equation}
\(\omega\le 0\) 或 \(\omega\ge 2\) 时 \(\rho(\mathcal L_{\omega})\ge 1\)。

\subsection{最佳松弛因子}

若 \(A\) 还具有相容次序，且 \(\rho(G_J)<1\)，则
\begin{equation}
\omega_{\mathrm{opt}}=\frac{2}{1+\sqrt{1-\rho(G_J)^{2}}}.
\end{equation}
\(\omega_{\mathrm{opt}}\in(1,2)\) 是使 \(\rho(\mathcal L_{\omega})\) 最小的松弛因子。模型问题上 \(\rho(G_J)=\cos(\pi h)\)，于是
\[
\omega_{\mathrm{opt}}=\frac{2}{1+\sin(\pi h)}.
\]

\subsection{渐近收敛速度}

在上述条件下
\begin{equation}
\rho(\mathcal L_{\omega_{\mathrm{opt}}})=\omega_{\mathrm{opt}}-1,
\end{equation}
故 \(R_{\infty}(\mathcal L_{\omega_{\mathrm{opt}}})=-\ln(\omega_{\mathrm{opt}}-1)\)。
网格加密时它仍趋于 \(0\)，但比 G--S 慢得轻得多：G--S 是 \(O(h^{2})\)，最佳 SOR 是 \(O(h)\)。

\subsection{超松弛理论的推广}

把一次扫描改成对称的两次（SSOR），或把 \(M\) 取成更一般的三角/块三角分裂，格式仍是单步线性定常 \(x^{(k+1)}=Gx^{(k)}+d\)，收敛仍由 \(\rho(G)<1\) 判定。
块 Jacobi、块 G--S、块 SOR 只是把 \(D,L,U\) 换成相应的块，分量公式换成块前代。

\section{习题}

\pending

\section{上机习题}

\pending

\chapter{共轭梯度法}

\section{最速下降法}

\pending

\section{共轭梯度法及其基本性质}

\subsection{共轭梯度法}

\pending

\subsection{基本性质}

\pending

\section{实用共轭梯度法及其收敛性}

\subsection{实用共轭梯度法}

\pending

\subsection{收敛性分析}

\pending

\section{预优共轭梯度法}

\pending

\section{Krylov 子空间法}

\subsection{正则化方法}

\pending

\subsection{残量极小化方法}

\pending

\subsection{残量正交化方法}

\pending

\section{习题}

\pending

\section{上机习题}

\pending

\chapter{非对称特征值问题的计算方法}

\section{基本概念与性质}

\pending

\section{幂法}

\pending

\section{反幂法}

\pending

\section{QR 方法}

\subsection{基本迭代与收敛性}

\pending

\subsection{实 Schur 标准形}

\pending

\subsection{上 Hessenberg 化}

\pending

\subsection{带原点位移的 QR 迭代}

\pending

\subsection{双重步位移的 QR 迭代}

\pending

\subsection{隐式 QR 算法}

\pending

\section{习题}

\pending

\section{上机习题}

\pending

\chapter{对称特征值问题的计算方法}

\section{基本性质}

\pending

\section{对称 QR 方法}

\subsection{三对角化}

\pending

\subsection{隐式对称 QR 迭代}

\pending

\subsection{隐式对称 QR 算法}

\pending

\section{Jacobi 方法}

\subsection{经典 Jacobi 方法}

\pending

\subsection{循环 Jacobi 方法及其变形}

\pending

\subsection{Jacobi 方法的并行方案}

\pending

\section{二分法}

\pending

\section{分而治之法}

\subsection{分割}

\pending

\subsection{胶合}

\pending

\section{奇异值分解的计算}

\subsection{二对角化}

\pending

\subsection{SVD 迭代}

\pending

\subsection{SVD 算法}

\pending

\section{习题}

\pending

\section{上机习题}

\pending


\end{document}
